Jedes große Unternehmen sammelt, generiert und verarbeitet riesige Mengen an Daten. Diese Daten stammen aus unterschiedlichsten Funktionen und Geschäftseinheiten eines Unternehmens und würden deshalb auch an verschiedenen Orten gespeichert werden, wenn es keine zentrale Lösung zur Verwaltung dieser Daten geben würde \cite[\pagef 123]{davenport_putting_1998}. 
Ein \ac{ERP}-System ist wörtlich ins Deutsche übersetzt ein \glqq Unternehmensressourcenplanungssystem\grqq{}, welches genau dieses Problem löst. Dabei handelt es sich um ein Standardsoftwarepaket, das Geschäftslogiken bereitstellt, die die Kernprozesse und Verwaltungsfunktionen eines Unternehmens unterstützen und diese zentral in einem System integrieren, wodurch eine hohe Funktionalität bereitgestellt werden kann. Zu diesen Geschäftsprozessen gehören beispielsweise typische Unternehmensfunktionen wie Beschaffung, Materialwirtschaft, Produktion, Logistik, Vertrieb oder Finanzbuchhaltung. Dabei fokussiert sich das \ac{ERP}-System auf die Integration wiederkehrender Geschäftsprozesse, wie Bestell- und Auftragsabwicklung oder Zahlungsprozesse, und lässt Prozesse, die weniger strukturiert oder unregelmäßig ablaufen, wie beispielsweise Marketing oder Produktentwicklung, eher außer Acht \cite[\pagef 143]{klaus_what_2000}.

\ac{ERP}-Systeme unterstützen die Ausführung von Geschäftsprozessen allerdings nicht nur isoliert innerhalb einzelner Funktionsbereiche, sondern in integrierter Form über mehrere Unternehmensbereiche hinweg. Dabei werden zusammenhängende Prozessschritte, wie etwa Auftragserfassung, Logistik und Finanzbuchhaltung, systemintern miteinander verknüpft und automatisch koordiniert. Auf diese Weise fungieren \ac{ERP}-Systeme als zentrale Grundlage für durchgängige Geschäftsprozesse innerhalb eines Unternehmens. \cite[\pagef 1ff]{wortmann_evolution_1998}

Charakteristisch für ein \ac{ERP}-System ist zudem die Nutzung einer gemeinsamen, zentralen Datenbasis, auf die alle integrierten Module zugreifen. Dadurch wird eine konsistente und unternehmensweit einheitliche Datenhaltung ermöglicht. \cite[\pagef 123]{davenport_putting_1998}

Zusammenfassend definieren Klaus et al. ein \ac{ERP}-System als eine umfassende Softwarelösung, die darauf abzielt, alle Prozesse und Funktionen eines Unternehmens zu integrieren und in einer einheitlichen Informations- und IT-Architektur abzubilden, um dadurch eine ganzheitliche Sicht auf das Unternehmen zu ermöglichen \cite[\pagef 141f]{klaus_what_2000}.